\documentclass[11pt]{article}

\usepackage[margin=1in]{geometry}
\usepackage{enumitem}
\usepackage{titlesec}
\usepackage{hyperref}
\usepackage{xcolor}
\usepackage{tabularx}
\usepackage{booktabs}
\usepackage{longtable}

\titleformat{\section}{\large\bfseries}{}{0em}{}[\titlerule]
\titleformat{\subsection}{\normalsize\bfseries}{}{0em}{}

\hypersetup{colorlinks=true, linkcolor=blue, urlcolor=blue}

\begin{document}

% ============================================================
% HEADER
% ============================================================
\begin{center}
    {\LARGE \textbf{Status Report: Sprint \#1}} \\[6pt]
    {\large CS161, Spring 2026}
\end{center}

\vspace{0.5em}

% ============================================================
% 1. Team Information
% ============================================================
\section{Team Information}

\begin{tabularx}{\textwidth}{@{}lX@{}}
    \textbf{Team Number:}          & Team 2 \\
    \textbf{Team Name:}            & Checkpoint \\
    \textbf{Team Members:}         & Omkar Podey, Anmol Rastogi \\
    \textbf{Class:}                & CS161, Software Engineering, Spring 2026 \\
    \textbf{Date:}                 & March 15, 2026 \\
    \textbf{Sprint Period:}        & February 22 to March 15, 2026 \\
\end{tabularx}

% ============================================================
% 2. Task List Link
% ============================================================
\section{Product Backlog / Task List}

\textbf{Trello Board:} \url{https://trello.com/b/ufZ3LwJi/checkpoint-cs161}

\textbf{GitHub Repository:} \url{https://github.com/Omkar399/checkpoint}

% ============================================================
% 3. Sprint 1 Tasks --- Completed vs. Incomplete
% ============================================================
\section{Sprint 1: Task Summary}

All 8 Sprint~1 tasks were \textbf{completed} by the end of the sprint. Tasks were evenly divided between the two team members (4 each).

\begin{center}
\begin{tabularx}{\textwidth}{@{}lXcc@{}}
\toprule
\textbf{\#} & \textbf{Task} & \textbf{Assigned To} & \textbf{Status} \\
\midrule
1 & User Authentication (Login/Register) & Omkar & \textbf{Done} \\
2 & Invite Code System & Omkar & \textbf{Done} \\
3 & Channel Management & Omkar & \textbf{Done} \\
4 & Database Schema \& API Validation & Omkar & \textbf{Done} \\
5 & Server Creation & Anmol & \textbf{Done} \\
6 & Unified Chat Feed & Anmol & \textbf{Done} \\
7 & WebSocket Real-Time Messaging & Anmol & \textbf{Done} \\
8 & Polling Fallback Implementation & Anmol & \textbf{Done} \\
\bottomrule
\end{tabularx}
\end{center}

\textbf{Incomplete / Carried Over:} None. All Sprint~1 user stories were completed.

% ============================================================
% 4. What Went Well
% ============================================================
\section{What Went Well}

\begin{enumerate}[leftmargin=1.5em]
    \item \textbf{Clear task division}:Splitting the work into backend-focused (Omkar) and frontend-focused (Anmol) tracks allowed parallel development with minimal merge conflicts.

    \item \textbf{Architecture-first approach}:Spending time on the PRD and technical architecture before coding paid off. Having a defined database schema, API contract, and component hierarchy meant both developers could work independently with confidence.

    \item \textbf{All 8 features completed on time}:Every Sprint~1 task was finished within the two-week sprint window, keeping the project on track for Sprint~2.

    \item \textbf{Real-time messaging working end-to-end}:The WebSocket implementation with polling fallback was a higher-risk feature, and getting it functional in Sprint~1 de-risks the remainder of the project.

    \item \textbf{Tech stack choices}:FastAPI + SQLAlchemy/SQLite for backend and React + Vite for frontend proved to be productive choices, enabling rapid iteration without infrastructure overhead.
\end{enumerate}

% ============================================================
% 5. Issues / What Didn't Go Well
% ============================================================
\section{Issues / What Didn't Go Well}

\begin{enumerate}[leftmargin=1.5em]
    \item \textbf{Frontend/backend integration testing was delayed}: While individual components were built and tested in isolation, end-to-end integration testing across the full stack was not performed until late in the sprint. This caused some last-minute debugging around CORS configuration and WebSocket token passing.

    \item \textbf{WebSocket reliability on campus Wi-Fi}:The university network occasionally drops WebSocket connections. The polling fallback mitigates this, but the user experience during reconnection could be smoother.

    \item \textbf{No automated tests written}:Due to time constraints, the team focused on feature implementation over writing unit or integration tests. This is technical debt that should be addressed.

    \item \textbf{Trello cards lacked detail}:Cards were created with titles and assignments only, with no descriptions, checklists, or due dates were filled in. This made it harder to track granular progress during the sprint.
\end{enumerate}

% ============================================================
% 6. Changes for Next Sprint
% ============================================================
\section{Changes for Next Sprint}

\begin{enumerate}[leftmargin=1.5em]
    \item \textbf{Add descriptions and checklists to Trello cards}:Break each feature into sub-tasks with checkboxes so progress is more visible and granular.

    \item \textbf{Integration testing earlier}:Test frontend/backend integration continuously rather than deferring it to the end of the sprint.

    \item \textbf{Add due dates to cards}:Set target completion dates for each task to better track pacing within the sprint.

    \item \textbf{Begin writing basic tests}:Add at least API endpoint tests for critical flows (auth, check-ins) to catch regressions as the codebase grows.
\end{enumerate}

% ============================================================
% 7. Scrum Meetings
% ============================================================
\section{Scrum Meetings}

A total of \textbf{14 scrum meetings} were held during Sprint~1 (February 22 to March 15, 2026). Meetings were kept to 5 to 8 minutes each, covering what was done, what's next, and any blockers.

\begin{center}
\begin{tabularx}{\textwidth}{@{}clXc@{}}
\toprule
\textbf{\#} & \textbf{Date} & \textbf{Location} & \textbf{Attendance} \\
\midrule
1  & Feb 22 (Sat) & Discord (virtual)  & 100\% \\
2  & Feb 24 (Mon) & In-class           & 100\% \\
3  & Feb 25 (Tue) & Discord (virtual)  & 100\% \\
4  & Feb 26 (Wed) & In-class           & 100\% \\
5  & Feb 28 (Fri) & Discord (virtual)  & 100\% \\
6  & Mar 1 (Sat)  & Discord (virtual)  & 100\% \\
7  & Mar 3 (Mon)  & In-class           & 100\% \\
8  & Mar 4 (Tue)  & Discord (virtual)  & 100\% \\
9  & Mar 5 (Wed)  & In-class           & 100\% \\
10 & Mar 7 (Fri)  & Discord (virtual)  & 100\% \\
11 & Mar 10 (Mon) & In-class           & 100\% \\
12 & Mar 11 (Tue) & Discord (virtual)  & 100\% \\
13 & Mar 12 (Wed) & In-class           & 100\% \\
14 & Mar 14 (Fri) & Discord (virtual)  & 100\% \\
\bottomrule
\end{tabularx}
\end{center}

\textbf{Meeting breakdown:}
\begin{itemize}[leftmargin=1.5em]
    \item \textbf{6 in-class meetings}: Used class time for sprint planning, task coordination, and integration discussions.
    \item \textbf{8 virtual meetings (Discord)}: Held outside of class for implementation progress check-ins, code review, and blocker resolution.
\end{itemize}

% ============================================================
% 8. Individual Contributions
% ============================================================
\section{Individual Contributions}

\subsection{Omkar Podey}
\begin{itemize}[leftmargin=1.5em]
    \item Designed the database schema (7 tables: users, servers, server\_members, channels, channel\_members, messages, invites).
    \item Implemented JWT-based authentication with bcrypt password hashing (register, login, token validation).
    \item Built the invite code system with expiration, usage limits, and server join flow.
    \item Implemented channel management with CRUD operations, membership tracking, and authorization checks.
    \item Set up Pydantic request/response schemas for API validation across all endpoints.
    \item Configured the FastAPI project structure with SQLAlchemy ORM and SQLite.
\end{itemize}

\subsection{Anmol Rastogi}
\begin{itemize}[leftmargin=1.5em]
    \item Built the server creation feature with owner/member role management.
    \item Developed the unified chat feed with message pagination, grouping, and auto-scroll.
    \item Implemented WebSocket real-time messaging with the ConnectionManager and broadcast system.
    \item Built the 10-second polling fallback for unreliable network conditions.
    \item Created the React frontend scaffolding with Vite, including routing and layout components.
    \item Implemented the API client layer, Auth context provider, and WebSocket hook.
\end{itemize}

\end{document}
